Le Projet ITInéraire que nous avons mené durant cette deuxième partie de semestre consistait à développer une application concurrente à Strava, dédiée à l'enregistrement des activités sportives par GPS. Pour cela nous avions à notre disposition différents capteurs et une Raspberry Pi.


L'équipe pour ce projet était composé de 5 membres qui se sont occupés des différents aspects selon leurs affinités. On pourra séparer le projet en 3 grandes parties : 
\begin{itemize}
\item Les capteurs et leur gestion;
\item Le traitement statistique de calibration et des données brutes;
\item L'application et ses fonctionnalités;
\end{itemize}
La répartition de cu travail de chaque membre est plus approfondie dans la conclusion avec la matrice répartition.


Le projet s'est déroulé en plusieurs étapes. Premièrement nous avons du réaliser des acquisitions avec les différents capteurs afin d'évaluer leurs performances et notamment leur précision et erreurs grâce à des etudes statistiques des données récupérées. L'étude de ces différentes données des capteurs nous a permis ensuite de mieux les utiliser et les intégrer dans notre projet. En parallèle nous avons développé une interface ainsi que différentes fonctionnalités à destination des utilisateurs permettant de rendre plus lisible les données et analyser les performances sportives.


Finalement, après correction des capteurs nous avons pu développer un algorithme permettant de traiter les données brutes avant la mise en commun de toutes les parties du projet dans une seule et même application.


Nous décrirons dans ce rapport les différents capteurs qui étaient à notre disposition, le système d'acquisition créé pour le projet ainsi que l'application développée. La dernière partie du rapport sera consacrée à l'analyse des performances du système.
